%Copyright 2016 R.D. Martin
%This book is free software: you can redistribute it and/or modify it under the terms of the GNU General Public License as published by the Free Software Foundation, either version 3 of the License, or (at your option) any later version.
%
%This book is distributed in the hope that it will be useful, but WITHOUT ANY WARRANTY; without even the implied warranty of MERCHANTABILITY or FITNESS FOR A PARTICULAR PURPOSE.  See the GNU General Public License for more details, http://www.gnu.org/licenses/.
\section{Problems}
\begin{problem}{Two masses connected by a spring}
Two masses $m_1$ and $m_2$ are connected by a spring with spring constant $k$.\\
\textbf{a)} Write the Lagrangian for this system and show that it can be reduced to a single degree of freedom\\
\textbf{b)} Plot the effective potential for the system, and show that all orbits are bound\\
\textbf{c)} Determine the shape of the orbit, and plot it.
\end{problem}

\begin{problem}{The Coulomb force}
Two masses $m_1$ and $m_2$ with charges $q_1$ and $q_2$ interact through the Coulomb force.\\
\textbf{a)} Write the Lagrangian for this system and show that it can be reduced to a single degree of freedom\\
\textbf{b)} Plot the effective potential for the system when the force is attractive and describe the possible orbits\\
\textbf{c)} Plot the effective potential for the system when the force is repulsive and show that only hyperbolae are possible orbits\\
\textbf{d)} Plot the effective potential for the system when the force is repulsive but is constrained to act within a radius $r<R$ (for example, the Coulomb force from a nucleus that is screened by electrons at a large distance) and describe the possible orbits. 
\end{problem}

\begin{problem}{Kepler's problem}
\textbf{a)} Show that the analytic treatment presented in this chapter is consistent with Kepler's three laws\\
\textbf{b)} Show that for elliptic orbits, the semi-major axis only depends on the energy of the system.\\
\textbf{c)} Make a plot showing how the effective potential depends on angular momentum (show several curves for different values of angular momentum and comment).
\end{problem}

\begin{problem}{Simple Harmonic Oscillator}
\textbf{a)} Discuss the solutions for the simple harmonic oscillator for the cases when $E$ and/or $k$ are negative. If the solution exists, plot the position as a function of time.
\end{problem}

\begin{problem}{Coupled oscillators}
\textbf{a)} Show that in the case illustrated in example \ref{ex:threecoupledoscillators}, the total momentum is conserved. In particular, show that the Lagrangian can be written with one less degree of freedom.
\end{problem}
